\subsection{Environments and the pooled
marginal}\label{environments-and-the-pooled-marginal}

A latent environment \(E\) takes values in \(\{1,\dots,K\}\) with
probabilities \(\pi_e>0\). The observed vector follows a fixed
deterministic mixing,

\begin{equation}\label{eq:obs-model}
X=A\,x^{(0)}+b,
\end{equation}

where \(A\) has full column rank and \(b\) is a fixed offset;
conditional on \(E\), \(X\) has law \(G_e\). Without labels the analyst
receives only

\begin{equation}\label{eq:pool-mixture}
P_X=\sum_{e=1}^K \pi_e G_e,
\end{equation}

with no mark of the component behind a draw. Its second moment,

\begin{equation}\label{eq:pool-cov}
\bar\Sigma=E[XX^\top]=\sum_e\pi_e\Sigma_e,
\end{equation}

is one fixed positive-definite matrix, not a random per-sample
covariance. A hierarchical model with random \(C\) is a different object
we do not treat.

\subsection{Observation mechanisms and the recording
hierarchy}\label{observation-mechanisms-and-the-recording-hierarchy}

Two mechanisms are kept separate, with no implication between them.
\textbf{Mechanism A (one pooled distribution).} The analyst observes
draws from \(P_X\) or its moments and asks whether a mixture
decomposition is recoverable. \textbf{Mechanism B (repeated
recordings).} The analyst observes \(M\) recordings, each with several
windows and a covariance, and asks whether recordings can be assigned to
their environment.

Under mechanism B, a recording \(R_r\) holds \(W_r\) windows and carries
one environment label \(Z_r\) shared by every window. Recordings are the
independent units, windows the repeated measurements. Each window
covariance is estimated by shrinkage and linearized as
\(x_w=\operatorname{vech}(\log\widehat\Sigma_w)\); the recording feature
is the within-recording average,

\begin{equation}\label{eq:rec-estimand}
x_r=\frac1{W_r}\sum_{w=1}^{W_r}x_w,
\end{equation}

followed by one standardization across recordings, which makes differing
window counts and durations comparable.

\subsection{The two-layer estimand}\label{the-two-layer-estimand}

The window-level mean \(\mu_{e,w}=E[x_w^{(0)}\mid Z_r=e]\) and the
recording-level mean \(\mu_{e,\mathrm{rec}}=E[x_r^{(0)}\mid Z_r=e]\)
must not be conflated. DELVE targets the recording-level object. The two
coincide only under extra conditions on the within-recording averaging.
The finite-window target \(\mu_{e,T}\) approaches the window-level
target as the effective window count grows.

The back-transformed summary
\(\widetilde\Sigma_e=\operatorname{expm}(\mu_{e,w})\) is a log-Euclidean
summary, not the component covariance \(\Sigma_e\). An \(O(1/T_w)\)
conditional bias, including the shrinkage contribution, relates the two.
The log-Euclidean mean is also distinct from the affine-invariant
Karcher mean. We never treat a cluster center as a component covariance.

\subsection{Identification objects, environments, and separation
quantities}\label{identification-objects-environments-and-separation-quantities}

Three objects carry distinct conclusion forms. \textbf{L1 (component
parameters).} Under mechanism A, the unordered component parameters are
identified up to permutation. Finite moments reach generic injectivity
only at a higher order. \textbf{L2 (recording assignment).} Under
mechanism B, the assignment is stated as a misclassification-rate or
posterior-consistency result rather than a mixture theorem. It needs a
constant within-recording environment, consistent features, and a
positive margin. \textbf{L3 (single-window classification).} The window
label carries a Bayes error from distributional overlap that
recording-level consistency does not remove. A benign label permutation
is distinguished from genuine non-identification, where components
merge, the margin vanishes, or distributions overlap.

All statements condition on \(K\). The theory does not prove \(K\)
identifiable or consistently estimable. Automatic selection is an
engineering rule over \(\mathcal K=\{2,\dots,6\}\) that combines BIC,
silhouette, and bootstrap stability terms. Oracle-\(K\) is reported
separately from estimated-\(K\). Selection is a data-driven choice
within a candidate range, never fully blind discovery.

Three separation quantities are not interchangeable: \(\sigma_J\), the
smallest singular value of the Jacobian or symmetric-basis operator;
\(\Delta_{\mathrm{cov}}=\min_{e\ne f}\|\mu_e-\mu_f\|\), the assignment
margin; and \(\Delta_{\mathrm{dist}}\), a KL or total-variation distance
governing posterior rates and the Bayes error. Component-parameter
identification is necessary but not sufficient for assignment.
Assignment also needs enough recordings per environment and a positive
\(\Delta_{\mathrm{cov}}\). These requirements are stated as premises,
not derived from one another.
